% ============================================================================
% Welcome Page — Grade 7 Math Study Guide
% Hexagon-motif header with layered teal bands, centred welcome prose,
% and a 2×2 left-bordered showcard grid. THE TPT thumbnail page.
% UNIQUE to Study Guide: hexagon motifs, teal brand colour, left-bar cards.
% ============================================================================

\thispagestyle{empty}

% ── TikZ full-page background ────────────────────────────────────────────
\begin{tikzpicture}[remember picture, overlay]
    % Soft warm wash
    \fill[funTeal!2]
        (current page.south west) rectangle (current page.north east);

    % ── Layered header — three soft overlapping bands ──
    % Outermost: wide teal band with angled bottom
    \fill[funTeal!7]
        (current page.north west) --
        (current page.north east) --
        ([yshift=-9.8cm]current page.north east) --
        ([yshift=-10.6cm]current page.north west) -- cycle;

    % Middle: slightly inset green tint for depth
    \fill[funGreen!5]
        ([xshift=0.6cm, yshift=-0.6cm]current page.north west) --
        ([xshift=-0.6cm, yshift=-0.6cm]current page.north east) --
        ([xshift=-0.6cm, yshift=-9.4cm]current page.north east) --
        ([xshift=0.6cm, yshift=-10.1cm]current page.north west) -- cycle;

    % Inner: lighter lift
    \fill[white, opacity=0.35]
        ([xshift=1.6cm, yshift=-1.6cm]current page.north west) --
        ([xshift=-1.6cm, yshift=-1.6cm]current page.north east) --
        ([xshift=-1.6cm, yshift=-8.8cm]current page.north east) --
        ([xshift=1.6cm, yshift=-9.4cm]current page.north west) -- cycle;

    % ── Hexagon cluster — top left ──
    \node[regular polygon, regular polygon sides=6, fill=funTeal!12,
          minimum size=8mm, inner sep=0pt]
        at ([xshift=2.4cm, yshift=-1.6cm]current page.north west) {};
    \node[regular polygon, regular polygon sides=6, fill=funGreen!8,
          minimum size=5mm, inner sep=0pt]
        at ([xshift=3.8cm, yshift=-2.7cm]current page.north west) {};
    \node[regular polygon, regular polygon sides=6, fill=funTeal!7,
          minimum size=3.5mm, inner sep=0pt]
        at ([xshift=5.4cm, yshift=-1.3cm]current page.north west) {};

    % ── Hexagon cluster — top right ──
    \node[regular polygon, regular polygon sides=6, fill=funGreen!9,
          minimum size=6.5mm, inner sep=0pt]
        at ([xshift=-2.6cm, yshift=-1.9cm]current page.north east) {};
    \node[regular polygon, regular polygon sides=6, fill=funTeal!10,
          minimum size=4mm, inner sep=0pt]
        at ([xshift=-4.4cm, yshift=-2.9cm]current page.north east) {};
    \node[regular polygon, regular polygon sides=6, fill=funTeal!6,
          minimum size=2.8mm, inner sep=0pt]
        at ([xshift=-6.2cm, yshift=-1.5cm]current page.north east) {};

    % ── Hexagon cluster — bottom corners ──
    \node[regular polygon, regular polygon sides=6, fill=funGreen!5,
          minimum size=5mm, inner sep=0pt]
        at ([xshift=2.6cm, yshift=2.3cm]current page.south west) {};
    \node[regular polygon, regular polygon sides=6, fill=funTeal!5,
          minimum size=3.5mm, inner sep=0pt]
        at ([xshift=5.2cm, yshift=1.7cm]current page.south west) {};
    \node[regular polygon, regular polygon sides=6, fill=funTeal!6,
          minimum size=4.5mm, inner sep=0pt]
        at ([xshift=-3cm, yshift=2.6cm]current page.south east) {};
    \node[regular polygon, regular polygon sides=6, fill=funGreen!4,
          minimum size=2.5mm, inner sep=0pt]
        at ([xshift=-5.4cm, yshift=1.5cm]current page.south east) {};

    % ── Tracking text ──
    \node[font=\fontsize{9}{11}\selectfont\sffamily\color{funTealDark!50}]
        at ([yshift=-2.4cm]current page.north)
        {\textls[180]{YOUR GRADE 7 STUDY GUIDE}};

    % ── Main title ──
    \node[font=\fontsize{34}{40}\selectfont\sffamily\bfseries\color{black!88},
          text width=0.82\paperwidth, align=center]
        at ([yshift=-4.6cm]current page.north)
        {Math Made Easy};

    % ── Subtitle ──
    \node[font=\normalsize\sffamily\color{black!45},
          text width=0.7\paperwidth, align=center]
        at ([yshift=-6.2cm]current page.north)
        {Key Concepts \;\textbullet\; Worked Examples \;\textbullet\;
        Review \& Practice};

    % ── Elegant separator: triple-hexagon motif ──
    \draw[funTeal!22, line width=0.9pt]
        ([xshift=3.5cm, yshift=-7.5cm]current page.north west) --
        ([xshift=-3.5cm, yshift=-7.5cm]current page.north east);
    % Central hexagon
    \node[regular polygon, regular polygon sides=6, fill=funTeal!15,
          draw=funTeal!28, line width=0.4pt, minimum size=5.5mm,
          inner sep=0pt]
        at ([yshift=-7.5cm]current page.north) {};
    % Flanking hexagons
    \node[regular polygon, regular polygon sides=6, fill=funGreen!8,
          minimum size=3.5mm, inner sep=0pt]
        at ([xshift=-1.2cm, yshift=-7.5cm]current page.north) {};
    \node[regular polygon, regular polygon sides=6, fill=funGreen!8,
          minimum size=3.5mm, inner sep=0pt]
        at ([xshift=1.2cm, yshift=-7.5cm]current page.north) {};

    % ── Thin outer frame ──
    \draw[black!6, line width=0.3pt]
        ([xshift=8mm, yshift=-8mm]current page.north west) rectangle
        ([xshift=-8mm, yshift=8mm]current page.south east);

    % ── Copyright footer ──
    \node[font=\sffamily\tiny\color{black!32}, anchor=south]
        at ([yshift=0.5cm]current page.south)
        {\textcopyright~2026 Dr.\ A.\ Nazari \;\textbullet\; ViewMath.com};

\end{tikzpicture}

% ── Welcome Text ──────────────────────────────────────────────────────────
\vspace*{7.6cm}

\begin{center}
\begin{minipage}{0.76\textwidth}
    \centering\sffamily\large\color{black!68}
    \setlength{\parskip}{6pt}

    This study guide breaks down every Grade~7 math topic into
    clear explanations, worked examples, and practice problems.
    Read the concept, follow the examples step by step,
    then test yourself. The answer key in the back lets you
    check your work and learn from every mistake.
\end{minipage}
\end{center}

% ── Feature Showcards — 2×2 grid ─────────────────────────────────────────
\vspace{7mm}

\begin{center}
\begin{minipage}{0.86\textwidth}

\noindent
\begin{minipage}[t]{0.48\textwidth}
\sgShowcard{funTeal}{\faBookOpen}{9 Chapters}{%
    From ratios and percents to probability ---
    every Grade~7 standard explained clearly.}
\end{minipage}%
\hfill
\begin{minipage}[t]{0.48\textwidth}
\sgShowcard{funGreen}{\faLightbulb}{Key Concepts First}{%
    Each topic starts with the core idea, so you always
    know \emph{what} you're learning and \emph{why}.}
\end{minipage}

\vspace{1mm}

\noindent
\begin{minipage}[t]{0.48\textwidth}
\sgShowcard{funOrange}{\faPen}{Worked Examples}{%
    Step-by-step solutions show you exactly how to solve
    each type of problem before you try on your own.}
\end{minipage}%
\hfill
\begin{minipage}[t]{0.48\textwidth}
\sgShowcard{funPurple}{\faCheckDouble}{Practice with Answers}{%
    Every section ends with practice problems.
    Full answer key included in the back.}
\end{minipage}

\end{minipage}
\end{center}

\clearpage
